% =============================================================================
% AMR-Vortrag: Beamer-Präsentation
% Quellen: beamer.md, amr_sprechernotizen.md, amr_handout.md
% Build:   pdflatex amr_vortrag.tex  (oder: make amr_vortrag.pdf)
% =============================================================================
\documentclass[aspectratio=169]{beamer}
\usetheme{metropolis}

% -- Pakete ------------------------------------------------------------------
\usepackage{amr-dashboard}
\usepackage{graphicx}
\usepackage{booktabs}
\usepackage{amsmath}
\usepackage[T1]{fontenc}
\usepackage[utf8]{inputenc}
\usepackage[ngerman]{babel}
\usepackage{pgfpages}

% Sprechernotizen aktivieren (auskommentieren fuer Handout/Druck):
% \setbeameroption{show notes on second screen=right}

% -- Metadaten ---------------------------------------------------------------
\title{Effizienz auf der letzten Meile}
\subtitle{Konzeption und Validierung eines Navigationssystems\\
          für einen autonomen Kleinladungsträger-Transporter mit ROS~2}
\author{[Autor]}
\date{[Datum]}
\institute{[Hochschule / Fachbereich]}

% =============================================================================
\begin{document}
% =============================================================================

% -- Titelfolie --------------------------------------------------------------
\maketitle

\note{%
  Begrüßung des Publikums. Betonen, dass die Präsentation den
  kompletten Datenfluss vom physischen Sensorwert bis zur strategischen
  Cloud-Entscheidung verfolgt. Ziel: Zeigen, wie komplexe
  Logistikaufgaben mit einer stark partitionierten Hardware-Architektur
  gelöst werden können.%
}

% -- Gliederung --------------------------------------------------------------
\begin{frame}{Gliederung}
  \tableofcontents
\end{frame}

% =============================================================================
\section{Motivation und Problemstellung}
% =============================================================================

\begin{frame}{Motivation und Problemstellung}

  \textbf{Ausgangslage: Intralogistik im Wandel}
  \begin{itemize}
    \item Industrie~4.0 erhöht die Anforderungen an flexible Materialflüsse.
    \item Der Transport von Kleinladungsträgern bindet Personal ohne
          direkte Wertschöpfung.
    \item Autonome mobile Roboter arbeiten ohne Leitlinien oder Reflektoren.
  \end{itemize}

  \vspace{0.5em}
  \textbf{Problemstellung}
  \begin{itemize}
    \item Industrielle AMR wie der MiR100 liegen bei mehr als
          $25\,000\,\mathrm{EUR}$ und sind für einfache KLT-Transporte
          oft überdimensioniert.
    \item Für kostengünstige Open-Source-Systeme fehlen belastbare
          Referenzarchitekturen.
    \item Die zentrale technische Herausforderung besteht in der Kopplung
          von echtzeitfähiger Antriebsregelung und autonomer Navigation
          auf kostengünstiger Hardware.
  \end{itemize}

  \note{%
    Industrielle AMR kosten oft über 25.000~EUR und sind für einfache
    Aufgaben überdimensioniert. Unser Modell: Prototyp für rund 513~EUR
    Materialkosten. Die Kostenreduktion erfordert eine strikte
    Systemtrennung -- harte Echtzeit für Motorregelung und
    rechenintensive SLAM-Navigation dürfen sich nicht gegenseitig
    blockieren.%
  }
\end{frame}

% =============================================================================
\section{Zielsetzung und Projektfragen}
% =============================================================================

\begin{frame}{Zielsetzung und Projektfragen}

  \textbf{Hauptziel}\\
  Entwicklung und Validierung eines ROS-2-basierten Navigationssystems
  für einen kostengünstigen AMR mit Gesamtkosten von rund
  $500\,\mathrm{EUR}$.

  \vspace{0.5em}
  \textbf{Projektfragen}
  \begin{itemize}
    \item \textbf{PF1 -- Echtzeitarchitektur:} Wie lässt sich auf einer
          Dual-Knoten-Architektur mit zwei ESP32-S3 eine echtzeitfähige
          Regelung bei gleichzeitiger Sensorerfassung mit micro-ROS
          realisieren?
    \item \textbf{PF2 -- Navigationsgenauigkeit:} Welchen Einfluss haben
          UMBmark-Kalibrierung, IMU-Fusion und hardwarenahe
          Sicherheitslogik auf die Navigationsgenauigkeit?
    \item \textbf{PF3 -- Visuelle Wahrnehmung und Docking:} Erreicht ein
          monokulares Kamerasystem mit Edge-Inferenz und Cloud-Semantik
          eine ausreichende Präzision und Kontextsensitivität?
  \end{itemize}

  \vspace{0.3em}
  \textbf{Methodik:} V-Modell nach VDI~2206

\end{frame}

% =============================================================================
\section{Systemarchitektur}
% =============================================================================

\begin{frame}[fragile]{Systemarchitektur}

  \textbf{Verteilte Drei-Ebenen-Architektur}

  \begin{columns}[T]
    \begin{column}{0.55\textwidth}
      {\small\ttfamily
      \begin{tabbing}
        xxxx\=xxxxxxxxxxxxxxxxx\=\kill
        Drive-Knoten ESP32 \> $\longleftrightarrow$ \> Raspberry Pi~5\\
        \> Core~1: PID 50\,Hz \> Bedien- und\\
        \> Core~0: /odom \> Leitstandsebene\\[4pt]
        Sensor-Knoten ESP32 \> $\longleftrightarrow$ \> Lokalisierung und\\
        \> Core~1: I2C-Sensorik \> Kartierung\\
        \> Core~0: /cliff \> Navigation\\
        \> \> Sicherheitslogik
      \end{tabbing}}
    \end{column}
    \begin{column}{0.40\textwidth}
      {\small\ttfamily
      \begin{tabbing}
        xxx\=\kill
        Intelligente Interaktion\\
        \> Hailo-8L Edge-Inferenz\\
        \> Gemini-Semantik
      \end{tabbing}}
    \end{column}
  \end{columns}

  \vspace{0.5em}
  \begin{itemize}
    \item \textbf{Kommunikation:} micro-ROS über UART (USB-CDC) mit
          Reliable QoS
    \item \textbf{Synchronisation:} FreeRTOS-Mutex zwischen Core~0 und
          Core~1 auf beiden ESP32-S3
  \end{itemize}

  \note{%
    Bildführung: Von links nach rechts durch das Diagramm führen.
    Kernelement: Datenfluss Sensorik $\rightarrow$ Bildverarbeitung
    (Hailo) $\rightarrow$ Strategie (Cloud~AI) $\rightarrow$ Handlung
    (ESP32/Servos). Zeitkritische Notfallreaktionen (Stopp vor Wand)
    bleiben lokal auf dem Pi~5, semantische Entscheidungen gehen an die
    Cloud.%
  }
\end{frame}

% =============================================================================
\section{Hardware-Plattform}
% =============================================================================

\begin{frame}{Hardware-Plattform}

  \textbf{Komponenten mit Gesamtkosten von rund $500\,\mathrm{EUR}$}

  \vspace{0.3em}
  {\small
  \begin{tabular}{@{}ll@{}}
    \toprule
    \textbf{Komponente} & \textbf{Funktion} \\
    \midrule
    Raspberry Pi~5           & Bedien- und Leitstandsebene \\
    Hailo-8L (PCIe)          & Edge-Inferenz für Objekterkennung \\
    Drive-Knoten (ESP32-S3)  & Fahrkern: PID-Regelung und Rad-Odometrie \\
    Sensor-Knoten (ESP32-S3) & Sensor- und Sicherheitsbasis: IMU, Batterie, Kanten \\
    Cytron MDD3A             & Dual-Motortreiber im PWM-Betrieb \\
    2\,$\times$\,JGA25-370   & Getriebemotoren mit Hall-Encoder \\
    RPLIDAR A1               & 2D-LiDAR für Lokalisierung und Kartierung \\
    IMX296                   & Global-Shutter-Kamera \\
    \bottomrule
  \end{tabular}}

  \vspace{0.5em}
  \textbf{Zentrale Parameter:}
  Raddurchmesser $65{,}67\,\mathrm{mm}$,\;
  Spurbreite $178\,\mathrm{mm}$,\;
  Encoder ${\sim}\,748$ Ticks/Umdrehung,\;
  $K_p = 0{,}4$,\; $K_i = 0{,}1$,\;
  Zielgeschwindigkeit $0{,}4\,\mathrm{m/s}$

  \note{%
    RPLIDAR~A1: 360°-Abtastung bei 7,6~Hz bis 12~m Reichweite -- Basis
    für SLAM. IMX296: optische Daten für ArUco-Docking und
    YOLOv8-Erkennung. Hall-Encoder: Quadraturmodus, ${\sim}\,748$~Ticks
    pro Radumdrehung. Hier entsteht die reine Rohdaten-Grundlage ohne
    Interpretation.%
  }
\end{frame}

% =============================================================================
\section{Software-Stack}
% =============================================================================

\begin{frame}{Software-Stack}

  \textbf{Firmware auf den ESP32-S3}
  \begin{itemize}
    \item Der Fahrkern führt die PID-Regelschleife mit $50\,\mathrm{Hz}$
          aus und publiziert \texttt{/odom}.
    \item Die Sensor- und Sicherheitsbasis erfasst IMU-, Batterie- und
          Kanten-Signale und publiziert \texttt{/imu} sowie \texttt{/cliff}.
    \item Ein Komplementärfilter fusioniert Gyro- und Encoder-Anteile
          für die Heading-Schätzung.
  \end{itemize}

  \vspace{0.5em}
  \textbf{ROS~2 auf dem Raspberry Pi~5}
  \begin{itemize}
    \item ROS~2 Humble läuft containerisiert in Docker
          (\texttt{ros:humble-ros-base}, \texttt{arm64}).
    \item Der micro-ROS Agent wird aus dem Quellcode gebaut.
    \item \texttt{network\_mode: host} und \texttt{privileged: true}
          sichern den Zugriff auf DDS, serielle Schnittstellen und Geräte.
  \end{itemize}

  \vspace{0.5em}
  \textbf{Zentrale ROS-2-Knoten:}
  \texttt{micro\_ros\_agent},
  \texttt{cliff\_safety\_node},
  \texttt{dashboard\_bridge},
  \texttt{slam\_toolbox},
  \texttt{nav2}

  \note{%
    ROS~2 auf dem Pi~5: High-Level-Navigationsrechner, Docker-Container.
    Nav2 und SLAM: 5~cm genaue Occupancy Grid Map. Hailo-8L: NPU ist
    zwingend -- YOLOv8 würde CPU des Pi~5 blockieren. PCIe-Beschleuniger
    verarbeitet Videostream latenzfrei.%
  }
\end{frame}

% =============================================================================
\section{Navigation und hybride Vision}
% =============================================================================

\begin{frame}{Navigation und hybride Vision}

  \textbf{TF-Baum}

  {\ttfamily\small
    map $\rightarrow$ odom $\rightarrow$ base\_link $\rightarrow$ laser\\
    \hspace*{12.5em}$\rightarrow$ camera\_link\\
    \hspace*{12.5em}$\rightarrow$ ultrasonic\_link
  }

  \vspace{0.5em}
  \textbf{Lokalisierung und Kartierung sowie Navigation}
  \begin{itemize}
    \item \texttt{slam\_toolbox}: Ceres-Solver, Loop Closure,
          Rasterauflösung $5\,\mathrm{cm}$.
    \item Nav2: Zielanfahrt mit Regulated Pure Pursuit.
    \item Costmaps: \texttt{StaticLayer}, \texttt{ObstacleLayer},
          \texttt{InflationLayer}.
    \item Recovery: Drehen, Rücksetzen und Warten.
  \end{itemize}

  \vspace{0.5em}
  \textbf{Hybride Vision}
  \begin{itemize}
    \item \texttt{host\_hailo\_runner}: Bounding Boxes, Inferenzlatenz
          ${\sim}\,34\,\mathrm{ms}$.
    \item \texttt{gemini\_semantic\_node}: kontextbezogene Szenenanalyse.
    \item Docking: ArUco-Marker, SolvePnP-Pose, P-Regler zur
          Feinpositionierung.
  \end{itemize}

  \note{%
    Schnittstelle: Das System sendet vorverarbeitete JSON-Pakete an
    Gemini API oder Claude AI. Keine rohen Videostreams in die Cloud,
    sondern nur semantische Metadaten. Beispiel: LiDAR meldet Hindernis
    (lokaler Stopp), Kamera+Hailo erkennen ``Mensch'', Cloud-KI
    entscheidet strategisch (Warnung per Speaker statt Umplanung).%
  }
\end{frame}

% =============================================================================
\section{Validierung und Ergebnisse}
% =============================================================================

\begin{frame}{Validierung und Ergebnisse}

  \textbf{Verifikation der Subsysteme}

  \vspace{0.3em}
  {\small
  \begin{tabular}{@{}lll@{}}
    \toprule
    \textbf{Test} & \textbf{Ergebnis} & \textbf{Kriterium} \\
    \midrule
    PID-Regelfrequenz    & $50\,\mathrm{Hz}$, Jitter $< 2\,\mathrm{ms}$ & erfüllt \\
    Kanten-Stopp-Latenz  & $< 50\,\mathrm{ms}$                          & erfüllt \\
    Hailo-8L-Inferenz    & ${\sim}\,34\,\mathrm{ms}$/Frame              & erfüllt \\
    Odometrie-Rate       & stabile $20\,\mathrm{Hz}$                    & erfüllt \\
    Geradeausfahrt-Drift & $1{,}5\,\mathrm{cm}$ mit IMU-Fusion          & erfüllt \\
    \bottomrule
  \end{tabular}}

  \vspace{0.5em}
  \textbf{Validierung der Navigation}

  \vspace{0.3em}
  {\small
  \begin{tabular}{@{}lll@{}}
    \toprule
    \textbf{Test} & \textbf{Ergebnis} & \textbf{Kriterium} \\
    \midrule
    ATE bei Kartierung   & $0{,}16\,\mathrm{m}$ & $< 0{,}20\,\mathrm{m}$ \\
    Positionsgenauigkeit & $6{,}4\,\mathrm{cm}$  & $< 10\,\mathrm{cm}$    \\
    Winkelgenauigkeit    & $4{,}2^{\circ}$        & $< 8{,}6^{\circ}$      \\
    \bottomrule
  \end{tabular}}

  \note{%
    UMBmark-Kalibrierung hat systematische Odometriefehler um Faktor
    10--20 reduziert. Präzise Mechanik-Messungen sind essenziell.
    micro-ROS: MTU 512~Bytes, Reliable QoS ermöglicht Fragmentierung bis
    2048~Bytes -- zwingend für große Odometrie-Nachrichten. Ausblick:
    EKF-Sensorfusion (IMU + Encoder) für bessere Schlupfkompensation.%
  }
\end{frame}

% -- Docking-Ergebnisse (eigene Folie wegen Platzbedarf) --------------------
\begin{frame}{Validierung -- Docking}

  \textbf{Validierung des Dockings}

  \vspace{0.5em}
  {\small
  \begin{tabular}{@{}lll@{}}
    \toprule
    \textbf{Test} & \textbf{Ergebnis} & \textbf{Kriterium} \\
    \midrule
    Lateraler Versatz   & $1{,}3\,\mathrm{cm}$             & $< 2\,\mathrm{cm}$  \\
    Orientierungsfehler & $2{,}8^{\circ}$                   & $< 5^{\circ}$        \\
    Erfolgsquote        & $80\,\%$ (8 von 10 Versuchen)     & erfüllt              \\
    \bottomrule
  \end{tabular}}

\end{frame}

% =============================================================================
\section{Kernergebnisse}
% =============================================================================

\begin{frame}{Kernergebnisse}

  \textbf{PF1 bestätigt}
  \begin{itemize}
    \item Trennung in Fahrkern und Sensor- und Sicherheitsbasis
          stabilisiert die Echtzeitfähigkeit.
    \item Serielle Kopplung über UART liefert deterministischere
          Laufzeiten als WLAN.
  \end{itemize}

  \vspace{0.4em}
  \textbf{PF2 bestätigt}
  \begin{itemize}
    \item UMBmark-Kalibrierung korrigiert Raddurchmesser von
          $65{,}00$ auf $65{,}67\,\mathrm{mm}$.
    \item Systematische Odometriefehler sinken um Faktor 10--20.
    \item Sicherheitslogik stoppt zuverlässig vor Navigation.
  \end{itemize}

  \vspace{0.4em}
  \textbf{PF3 bedingt bestätigt}
  \begin{itemize}
    \item Hybride Vision erreicht beim Docking $80\,\%$ Erfolgsquote.
    \item Edge-Inferenz schnell genug für laufende Eingriffe.
    \item Grenzen bei Beleuchtung, Sichtfeld und Marker-Sichtbarkeit.
  \end{itemize}

  \vspace{0.3em}
  \textbf{Status:} Alle Muss-Anforderungen sind erfüllt.

\end{frame}

% =============================================================================
\section{Live-Demo}
% =============================================================================

\begin{frame}[fragile]{Live-Demo}

  \textbf{Demonstrationsszenario: KLT-Transport}
  \begin{enumerate}
    \item Lokalisierung und Kartierung der Umgebung
    \item Autonome Punkt-zu-Punkt-Navigation mit Objekterkennung
    \item Reaktion der Sicherheitslogik auf eine simulierte Kante
    \item ArUco-basiertes Docking an einer Ladestation
  \end{enumerate}

  \vspace{0.5em}
  \textbf{Systemstart}

  {\small\ttfamily
  ./run.sh ros2 launch my\_bot full\_stack.launch.py \textbackslash\\
  \quad use\_dashboard:=True use\_vision:=True use\_audio:=True
  }

  \vspace{0.5em}
  \textbf{Beobachtbare Topics}
  \begin{itemize}
    \item \texttt{/odom} -- Rad-Odometrie mit $20\,\mathrm{Hz}$
    \item \texttt{/cliff} -- Kanten-Erkennung mit $20\,\mathrm{Hz}$
    \item \texttt{/map} -- Belegungskarte aus \texttt{slam\_toolbox}
  \end{itemize}

\end{frame}

% =============================================================================
\section{Zusammenfassung und Ausblick}
% =============================================================================

\begin{frame}{Zusammenfassung und Ausblick}

  \textbf{Zusammenfassung}
  \begin{itemize}
    \item Vollständiger Entwicklungszyklus nach VDI~2206
    \item Verteilte Drei-Ebenen-Architektur mit zwei ESP32-S3 und einem
          Raspberry Pi~5
    \item Positive Beantwortung aller drei Projektfragen mit
          abgestuftem Ergebnis für PF3
    \item Erfüllung aller Muss-Anforderungen des Lastenhefts
  \end{itemize}

  \vspace{0.5em}
  \textbf{Ausblick}
  \begin{itemize}
    \item \textbf{Hardware:} Eigene Leiterplatte statt freier
          Verdrahtung, robustes Gehäuse
    \item \textbf{Sensorfusion:} EKF-Integration mit
          \texttt{robot\_localization}
    \item \textbf{Flottenmanagement:} Einbindung in Open-RMF
    \item \textbf{Sicherheit:} Weitere Annäherung an ISO~3691-4 und
          CE-Konformität
  \end{itemize}

\end{frame}

% =============================================================================
% Schlussfolie
% =============================================================================

\begin{frame}[standout]
  Fragen?
\end{frame}

% =============================================================================
% Appendix: Handout-Material
% =============================================================================

\appendix

\begin{frame}{Appendix -- Kinematik und Regelung}

  Inverskinematik des Differentialantriebs:
  \begin{align}
    \omega_l &= \frac{v - \omega \cdot \tfrac{L}{2}}{r} \\[4pt]
    \omega_r &= \frac{v + \omega \cdot \tfrac{L}{2}}{r}
  \end{align}

  \textbf{Systemparameter:}
  \begin{itemize}
    \item Zielgeschwindigkeit ($v$): $0{,}4\,\mathrm{m/s}$
          (Regulated Pure Pursuit Controller)
    \item Spurbreite ($L$): $178\,\mathrm{mm}$
    \item Radradius ($r$): $32{,}835\,\mathrm{mm}$
          (kalibrierter Raddurchmesser $65{,}67\,\mathrm{mm}$)
  \end{itemize}

  \note{%
    Herleitung über Inverskinematik. Ohne exakte, empirisch ermittelte
    Werte produziert der Nav2-Stack fehlerhafte Befehle, da Translation
    und Rotation physikalisch falsch auf die Räder abgebildet werden.%
  }
\end{frame}

\begin{frame}{Appendix -- Randbedingungen und Limitierungen}

  \begin{itemize}
    \item \textbf{Kommunikationsbandbreite:}
          Serielle UART-Verbindung ($921\,600\,\mathrm{Baud}$) nutzt
          XRCE-DDS-Protokoll. MTU auf $512\,\mathrm{Bytes}$ limitiert.
          Fragmentierung über Reliable QoS zwingend erforderlich für
          größere Pakete (z.\,B.\ Odometrie-Kovarianzmatrizen).

    \item \textbf{Cloud-Abhängigkeit:}
          Die strategische Objektauswertung setzt eine lückenlose
          Netzwerkverbindung voraus. Bei Verbindungsabbruch greift das
          System auf die rein lokale, reaktive Hindernisvermeidung des
          Nav2-Stacks zurück.
  \end{itemize}

\end{frame}

\begin{frame}{Appendix -- Datenfluss-Ebenen (Handout)}

  {\small
  \begin{enumerate}
    \item \textbf{Sensorik (Datenerfassung):}
          RPLIDAR A1 ($7{,}6\,\mathrm{Hz}$, $12\,\mathrm{m}$),
          IMX296 Kamera ($1456 \times 1088$\,px, $15\,\mathrm{fps}$),
          Hall-Encoder (${\sim}\,748\,\mathrm{Ticks/Umdrehung}$)

    \item \textbf{Edge-Computing (Lokale Verarbeitung):}
          Raspberry Pi~5 mit ROS~2 (Nav2, SLAM Toolbox),
          Hailo-8L NPU für YOLOv8-Objekterkennung

    \item \textbf{Cloud-KI (Strategische Entscheidung):}
          Gemini API / Claude AI empfängt aggregierte, semantische
          JSON-Daten

    \item \textbf{Aktorik (Physische Ausführung):}
          XIAO ESP32-S3 -- Core~1: PID bei $50\,\mathrm{Hz}$,
          Core~0: micro-ROS-Kommunikation.
          Cytron MDD3A: $20\,\mathrm{kHz}$ PWM
  \end{enumerate}}

\end{frame}

% =============================================================================
\end{document}
